% frontmatter-2C-v3-20261124.tex — main paper A (2C + critical scaling) frontmatter.
% NEW FILE, 2026-11-24. Supersedes frontmatter-2C-v2.tex for the assembled draft;
% v2 is left untouched (only-add discipline).
% Changes vs v2: the placeholder "[Critical scaling theorem: to be inserted]" is
% replaced by the proved/conditional statement; the honesty paragraph records that
% the pole decomposition is proved but the scaling theorem is conditional and
% neither is Lean-formalized; the companion sentence drops the page count.
\title{On window constraints, digit residues, and critical phase transitions
for modular continued fractions}

\begin{document}
\maketitle

\begin{abstract}
We determine Hausdorff dimensions of Hermite window-count
constraints: prescribed parity has exact dimension $1$,
and bounded or prescribed increments have exact dimension $1/2$; short windows
occur with explicit logarithmic frequency almost everywhere (constant
$(1-\ee^{-1})/160$), while divergent window counts force Hausdorff dimension
$1$ and Lebesgue nullity. For digit residues modulo $m\ge2$, eventual
periodicity has the same Hausdorff dimension as the fixed alphabet; a
uniform Lipschitz sandwich, an unconditional transfer-operator
criterion bypassing Bowen pressure, and exact
integer certificates yield the certified $m\in\{2,3,4,5,6\}$ brackets, the
global strict separation $\dimH F_{m,m}<\dimH F_{m,1}=\dimH E_m$ for all
$m\ge2$, and two-term expansion
$\dimH F_{m,r}=\frac12+\frac{C_r}{2m}-\frac{C_r^2\log m}{2m^2}+O(m^{-2})$
with $C_r$ the reciprocal metallic-continuant constants
($C_1=\psi_{\mathrm{Fib}}$). For multi-residue alphabets
the Bowen-free Poincar\'e-abscissa dimension is $s_0(S)$,
$\dimH F_{m,S,B}\to\max(1/2,s_0(S))$ with three-term Euler--Mascheroni
subcritical expansion to $O((\log m)^2/m^3)$ and supercritical rate
$\Theta(m^{-2s_0(S)})$. At the critical boundary $s_0(S)=1/2$ we prove
$C_S(p)=\kappa(S)/(p-2s_0)+O(1)$ with
$\kappa(S)=c(S)/(-\rho'(2s_0))$, and the critical exponent $1/2$ with amplitude
$\sqrt{d\kappa(S)}/2$ and no log correction.
All results are machine-checked in Lean~4 using only the base axioms
\texttt{propext}, \texttt{Classical.choice} and \texttt{Quot.sound}, except
the open-interval strictness of the certified brackets, the critical scaling
theorem, and the exact supercritical constant, which remain outside
the Lean kernel.
This work does not resolve the open problem of the irrationality of $e+\pi$.
\end{abstract}

\section*{Honesty statement and formal verification}

The Lean~4 audit subset of this paper --- the modules covering every theorem
stated above, with names transcribed verbatim from the corresponding sections
of the full manuscript --- is: the thirteen window modules
\texttt{WindowMobiusCount.lean}, \texttt{NoShortDimCount.lean},
\texttt{WindowTreePrune.lean}, \texttt{WindowAsympThresh.lean},
\texttt{WindowFrostmanDischarge.lean}, \texttt{WindowPaperScale.lean},
\texttt{WindowCFRealize.lean}, \texttt{WindowCFBridge.lean},
\texttt{WindowNoShortCFRealize.lean}, \texttt{WindowPrescribedCFRealize.lean},
\texttt{WindowBoundedUpper.lean}, \texttt{WindowShortAe.lean}, and
\texttt{WindowShortAeRealize.lean}; the residue modules
\texttt{EpiFormal/ResidueModGeneral.lean},
\texttt{EpiFormal/TransferOperatorDim.lean}, \texttt{EpiFormal/ResidueAsymptotic.lean}, and
\texttt{WindowAndResidueCore.lean}; the certificate modules
\texttt{EpiFormal/ResidueCertGenSound.lean},
\texttt{EpiFormal/ResidueCertGenBase.lean}, and
\texttt{EpiFormal/ResidueCertTable.lean}; and
\texttt{EpiFormal/ResidueMultiAsymptotic.lean}, audited separately in
\texttt{AuditResidueMultiAsymptotic.lean}. The full manuscript reports these
modules as kernel-checked at Layer~A with only the three base axioms and zero
\texttt{sorryAx}; they are audited jointly with the companion papers' modules
in \texttt{Audit.lean}, whose declaration counts are project-wide, not
per-paper.

The critical-scaling section of this paper is \emph{not} Lean-formalized. Its
pole decomposition is proved in the text by the Perron--Frobenius theorem for
the transfer operator on $A(\{|z-1|<1\})$ and by Kato's analytic perturbation
theory. The scaling theorem and the critical exponent are proved from the
master equation (Lemma \texttt{lem:master}), whose lower side uses the
finite-subalphabet transfer-operator criterion; that criterion's
step-certificate soundness is kernel-checked, but the analytic estimates of
the critical-scaling section are not Lean-formalized. The supercritical leading constant is a conjecture
supported numerically, not a theorem. Outside the Lean kernel also remain the
strictness of the open intervals of the narrow $m\in\{2,3,4,5,6\}$ brackets
(certified by the exact-integer Python script
\texttt{residue\_certificate\_general\_N512.py} together with the cited
pressure formalism; the closed brackets themselves are kernel-checked). The
collocation values for a possible fourth-order term are heuristic and are not
proved.

\emph{The full version containing all proofs is available as a companion manuscript.}
